
\chapter{EMV}\label{ch:emv}

\section{Магнитная полоса и чип}\label{sec:emv-magstripe-to-chip}

Магнитная полоса из~раздела~\ref{sec:card-data-magstripe} хранит
статические данные. PAN, срок действия, CVV (раскрыты в~разделах~\ref{sec:card-data-magstripe}
и~\ref{sec:card-data-cvv}) и~сервисный код записаны на~дорожке раз
и~навсегда при выпуске карты. Каждый считыватель получает одну и~ту
же последовательность байтов, и~каждая считанная копия неотличима
от~оригинала. Это свойство сделало магнитную полосу удобной мишенью
для мошенников. Достаточно один раз пропустить карту через скиммер,
накладку на~терминал или банкомат, незаметно считывающую дорожку,
и~полученные данные можно записать на~чистую карту-болванку, создав
полноценный клон.

В~1990-х платёжная отрасль искала способ снизить ущерб от~поддельных
карт через микрочип, который труднее скопировать, чем магнитную полосу.
Спецификация EMV (исходно аббревиатура от~Europay, Mastercard, Visa;
сегодня~--- имя стандарта, поддерживаемого EMVCo) версии~96 стала общей
платформой для такого перехода, а~в~1999~году появилась EMVCo, которая
взяла на~себя развитие и~сопровождение набора спецификаций. Сегодня
EMVCo принадлежит шести глобальным платёжным сетям~--- American Express,
Discover, JCB, Mastercard, UnionPay и~Visa~\cite{emvco-why-emv}; актуальная
редакция базовых контактных книг (Books~1--3)~--- v4.4 от~октября~2022~года,
Book~4 остаётся в~v4.3 от~ноября~2011~года~\cite{emvco-book1,emvco-book4}.

\highlight{В~EMV карта подтверждает подлинность вычислением,
а~не предъявлением статических данных}. Вместо статических данных,
которые терминал транслирует дальше как есть, чиповая карта
вычисляет криптограмму. Это значение зависит от~суммы
транзакции, валюты, даты, непредсказуемого числа от~терминала
(unpredictable number~--- случайное число, формируемое терминалом
для каждой транзакции)
и~секретного ключа, который никогда не покидает чип.
Скопировать криптограмму бесполезно: она валидна
только для одной транзакции, а~для следующей нужна новая,
вычислить которую без ключа невозможно~\cite{emvco-book1}.

Карточные сети дополнили спецификацию экономическим стимулом~---
переносом ответственности (\textit{liability shift};
полный разбор механики и~условий применения~--- в~разделе~\ref{sec:3ds-liability}):
риск мошенничества смещается на~ту сторону, которая не поддержала более
защищённый сценарий по~правилам конкретной сети. Visa для США
привязывает этот поворот к~1~октября 2015~года. Если чиповая карта
используется в~магазине, а~терминал не поддерживает чтение чипа,
риск ложится на~сторону, оставшуюся на~менее защищённой
технологии~\cite{visa-emv-liability-shift}.

Угроза немедленных чарджбэков и~отсроченных штрафов обновила парк
терминалов быстрее, чем это сделал~бы прямой регуляторный мандат.
Переход шёл волнами~--- сначала крупные рынки,
затем рынки с~более инерционной инфраструктурой.

EMV закрывает один класс мошенничества~--- клонирование
статических карточных данных в~операциях с~физическим предъявлением
карты. По~данным Visa, к~декабрю~2018 мошенничество с~поддельными
картами (\textit{counterfeit fraud}) в~канале с~физическим
предъявлением у~мерчантов США, перешедших на~чиповый приём,
упало на~80\,\% по~сравнению с~до-EMV уровнем 2015~года;
к~марту~2019~--- на~87\,\% у~модернизированных POS и~на~62\,\%
в~целом по~стране, считая и~мерчантов без
EMV~\cite{visa-emv-fraud-2019,american-banker-emv-2019}. Параллельно
фрод перетёк в~канал без предъявления карты (CNP, card not present),
где статическая защита остаётся прежней~--- отсюда экономическая
мотивация для 3-D~Secure (см.~раздел~\ref{sec:3ds-liability}).

\section{Архитектура чиповой транзакции}\label{sec:emv-architecture}

Чиповая карта, ICC (Integrated Circuit Card)~--- небольшая
вычислительная система со своим процессором,
криптографическим сопроцессором, оперативной памятью, постоянной
памятью и~перезаписываемой памятью; в~последней лежат данные эмитента,
счётчики транзакций и~ключи~\cite{emvco-book1}.


\begin{figure}[tbp]
  \centering
  \includefiguresvg[width=\textwidth]{assets/figures/ch07-emv/emv-chip-architecture}
  \caption{Архитектура чипа платёжной карты.}\label{fig:emv-chip-architecture}
  \medskip
\end{figure}

\subsection*{Внутри чипа: апплет и аттестация}\label{sec:emv-chip-internals}

В~ROM чипа лежит операционная система; для~платёжных карт это
Java~Card Classic. Платёжное приложение реализовано как \textit{апплет}
(applet), загружаемый в~перезаписываемую память при персонализации
и~выбираемый командой \texttt{SELECT} (см.~раздел~\ref{sec:emv-app-selection}). Управление
апплетами на~карте~--- загрузка, активация, блокировка, обновление
параметров~--- идёт через GlobalPlatform (отраслевой стандарт
управления контентом смарт-карт) по~защищённым каналам SCP02/SCP03
(Secure Channel Protocol, шифрованный и~аутентифицированный канал
поверх APDU); по~ним же эмитент доставляет issuer scripts (команды
эмитента к~чипу, разбор в~разделе~\ref{sec:emv-cryptogram})~\cite{globalplatform-card}.

Аппаратная база сертифицируется отдельно от~апплета. EMVCo ведёт
реестры одобренных интегральных микросхем (ICCN, Integrated~Circuit
Certificate Number) и~платформ (PCN, Platform Certificate Number);
с~2017~года процесс Chip~\& Platform Approval аккредитован
по~ISO/IEC~17065~\cite{emvco-cpa}.

Уровень безопасности самого чипа задаётся профилем защиты
Security~IC~Platform Protection~Profile (BSI-CC-PP-0084) по~схеме
Common~Criteria~--- международной системе оценки безопасности
ИТ-продуктов согласно ISO/IEC~15408~\cite{cc-pp-0084}. CC использует
две связанные шкалы~--- EAL (Evaluation Assurance Level, глубина
проверки) и~AVA\_VAN (Vulnerability ANalysis, стойкость к~атакам,
шкала 1--5). Базовая планка для платёжного чипа~--- EAL\,4+
с~AVA\_VAN.5 (атакующий с~неограниченным временем и~лабораторией);
промышленные NXP~JCOP, Infineon~SLE и~ST31 проходят на~EAL\,6+
с~тем же AVA\_VAN.5~\cite{nxp-jcop4-emv}. Российский аналог~---
ГОСТ~Р ИСО/МЭК~15408 с~оценочными уровнями доверия ОУД (ОУД4
соответствует EAL\,4).

Сертификация задаёт верхнюю планку атаки, которую чип обязан выдержать
без раскрытия ключа. Защиты на~\autoref{fig:emv-chip-architecture}:

\begin{itemize}
  \item TRNG (true random number generator, истинный ГСЧ~--- генератор
    случайных чисел) даёт свежие nonce;
  \item защитный экран поверх кристалла препятствует зондированию;
  \item сенсоры напряжения и~температуры регистрируют попытки внешнего
    воздействия;
  \item счётчики неудачных попыток необратимо стирают секреты
    при срабатывании сенсоров.
\end{itemize}

Они закрывают SPA/DPA (simple/differential power
analysis, анализ потребления тока микросхемой под нагрузкой),
fault~injection (glitch~--- кратковременный провал напряжения,
лазерный импульс, перегрев) и~физическое препарирование чипа.

Терминал общается с~чипом по~правилам ISO~7816. Этот стандарт
задаёт физический интерфейс с~контактами C1--C8 на~площадке, протокол
передачи данных T=0 или T=1 и~формат команд APDU
(Application Protocol Data Unit). Именно через APDU терминал читает
карту, выбирает приложение и~запрашивает обработку транзакции.
Команда APDU состоит из~заголовка с~классом, кодом инструкции и~двумя
параметрами и~необязательных данных; в~ответ карта возвращает данные
и~двухбайтовый код статуса SW1~SW2 (Status Word~1/2). Например, \texttt{0x9000} означает
успех, \texttt{0x6A82}~--- файл не найден, \texttt{0x6985}~--- условия
использования не выполнены. Первоисточник этих кодов~--- ISO/IEC~7816-4,
а~EMV Book~3 наследует их и~уточняет применение для платёжных
приложений~\cite{iso-7816-4,emvco-book3}.

Чиповая карта хранит несколько независимых платёжных приложений.
На~одной физической карте могут сосуществовать, например, приложение
Мир и~приложение UnionPay (кобейдж, выпускаемый в~России с~2017~года;
маршрутизация и~правила~--- в~главе~\ref{ch:mir})~\cite{nspk-jcb-cobadge},
каждое со своим набором ключей, сертификатов и~параметров. Какое именно
приложение пойдёт в~работу, определяется на~первом же шаге
сценария EMV~--- выборе приложения.

Весь контактный сценарий EMV состоит из~пяти последовательных фаз.

\begin{enumerate}
  \item \textbf{Выбор приложения}~--- терминал и~карта согласовывают,
    какое платёжное приложение будет использовано.
  \item \textbf{Аутентификация данных}~--- терминал проверяет,
    что данные на~карте подлинны и~не были подменены.
  \item \textbf{Верификация держателя}~--- карта или терминал
    проверяют, что операцию инициировал законный держатель,
    через PIN, подпись или иной метод.
  \item \textbf{Риск-менеджмент}~--- терминал и~карта выбирают между
    офлайн-одобрением и~онлайн-авторизацией у~эмитента.
  \item \textbf{Генерация криптограммы}~--- карта вычисляет
    криптографическое значение, подтверждающее результат всех
    предыдущих проверок.
\end{enumerate}

К~моменту генерации криптограммы карте доступны данные о~том,
какое приложение выбрано, прошла~ли аутентификация данных,
как завершилась проверка держателя и~нужно~ли обращаться
к~эмитенту онлайн~--- криптограмма фиксирует весь этот путь.

\section{Выбор приложения}\label{sec:emv-app-selection}

Прежде чем начать транзакцию, терминал должен определить, какое
платёжное приложение на~карте он будет использовать. Каждое такое
приложение идентифицируется \textit{AID} (Application Identifier)~---
по~нему терминал выбирает нужное ядро, набор правил и~сценарий
обработки. AID состоит из~двух частей: RID (Registered Application Provider
Identifier) карточной сети и~PIX (Proprietary Application Identifier
Extension) конкретного продукта. EMV~Book~1 описывает выбор в~два шага:
сначала терминал запрашивает у~карты каталог приложений, затем
согласовывает приложение из~списка кандидатов~\cite{emvco-book1}.

Терминал начинает с~обращения к~каталогу PSE
(Payment System Environment), а~в~бесконтактном сценарии~--- к~PPSE
(Proximity Payment System Environment),
описанному в~разделе~\ref{sec:contactless-emv}. В~ответ карта
возвращает список записей с~установленными приложениями, каждая
содержит AID. Из~этих записей терминал собирает
\emph{список кандидатов}~--- приложения, которые поддерживают и~карта,
и~терминал.

На~уровне протокола терминал посылает команду \texttt{SELECT}
к~файлу \texttt{1PAY.SYS.DDF01} или \texttt{2PAY.SYS.DDF01},
получает список записей с~короткими идентификаторами файлов
и~номерами записей и~извлекает из~них идентификаторы приложений.

Если PSE на~карте отсутствует~--- так бывает у~старых карт или
в~конфигурациях некоторых эмитентов,~--- терминал переходит
к~альтернативе. Он последовательно вызывает SELECT для каждого AID
из~своего списка поддерживаемых
и~смотрит, на~какой карта ответит положительно.

Когда в~списке кандидатов оказывается несколько приложений, терминал
сортирует их по~приоритету~--- тег \texttt{0x87}, Application Priority Indicator~---
и~либо автоматически выбирает приложение с~наивысшим приоритетом,
либо предлагает выбор держателю. В~большинстве торговых точек
выбор происходит автоматически, и~держатель видит на~экране терминала
лишь название выбранного приложения или запрос подтвердить выбор,
если того требует сценарий. У~Visa Debit/Credit (Classic) в~актуальном
TADG (Visa Technology Asset Documentation for Global Card Acceptance,
публичный технический справочник Visa) указан RID \texttt{0xA000000003},
PIX \texttt{0x1010} и~полный AID \texttt{0xA0000000031010}~\cite{visa-tadg}.

После успешного выбора приложения терминал отправляет команду
GPO (Get Processing Options). Она объявляет начало транзакции и~передаёт
первичные данные~--- сумму, валюту, дату и~другие параметры. Карта
в~ответ возвращает AIP (Application Interchange Profile, тег \texttt{0x82}), битовое
поле о~поддержке функций SDA, DDA, CDA (статической, динамической
и~комбинированной аутентификации данных; подробно~---
в~разделе~\ref{sec:emv-data-auth}), риск-менеджмента карты
и~верификации эмитента, и~AFL (Application File Locator), список
файлов и~записей, которые терминал должен прочитать для продолжения
транзакции~\cite{emvco-book3}. На~этом этапе терминал ничего ещё
не решает; он лишь выясняет, на~что карта способна и~какие данные
читать дальше.

Раскладка AIP по~битам приведена на~рис.~\ref{fig:aip-anatomy}:
перечисленные выше флаги уложены в~байт~1, большая часть байта~2
зарезервирована под бесконтактные расширения (RFU~CL, reserved
for EMV Contactless Specifications). В~разобранном примере
\texttt{AIP = 0x5800} карта поддерживает SDA, CVM и~требует
терминальный риск-менеджмент~\cite{emvco-book3}.

\begin{figure}[tbp]
  \centering
  \includefiguresvg[width=\linewidth]{assets/figures/ch07-emv/aip-anatomy}
  \caption{Раскладка AIP: 2~байта, MSB слева; жёлтым~--- биты примера \texttt{AIP = 0x5800}.}\label{fig:aip-anatomy}
\end{figure}

Какие именно параметры передавать в~GPO, диктует сама карта.
В~ответе на~SELECT она объявляет PDOL (Processing Options Data
Object List, тег \texttt{0x9F38})~--- список тегов и~их длин,
которые терминал должен сложить и~передать в~поле данных GPO.
Скажем, бесконтактная Visa объявляет PDOL вида
\begin{center}
  \small\texttt{0x9F66049F02069F37045F2A029F1A02}
\end{center}
\smallskip
\noindent (см.~рис.~\ref{fig:pdol-anatomy})~--- пять пар \texttt{(тег, длина)}: TTQ, сумма, непредсказуемое число, код валюты, код страны терминала. Терминал конкатенирует значения в~указанном порядке и~оборачивает их в~тег \texttt{0x83} в~теле GPO~\cite{emvco-book3,visa-tadg}.

Каждый тег занимает 2~байта (BER-TLV multi-byte indicator в~старших
разрядах первого байта), длина~--- 1~байт; на~рисунке байт длины
выделен более насыщенным оттенком.

\begin{figure}[tbp]
  \centering
  \includefiguresvg[width=.95\linewidth]{assets/figures/ch07-emv/pdol-anatomy}
  \caption{Побайтовый разбор 15-байтового PDOL бесконтактной Visa: пять пар \texttt{(тег, длина)}.}\label{fig:pdol-anatomy}
\end{figure}

Если PDOL у~карты нет~--- так бывает у~контактных Mastercard,~---
GPO уходит с~пустым телом \texttt{0x8300}. Реализации POS (Point of
Sale, торговый терминал), ожидавшие PDOL у~любой карты без~исключений,
на~таких профилях падали.

\subsection*{BER-TLV в~EMV: минимальный парсер}

PDOL~--- лишь один из~примеров TLV-структуры (Tag-Length-Value),
причём усечённый: в~нём нет значений, только пары
\texttt{(тег, длина)}. Полноценный BER-TLV (Basic Encoding Rules TLV)~---
цепочки \texttt{(тег, длина, значение)}~--- лежит
в~ответах карты на~SELECT, READ RECORD и~в~содержимом DE~55
(Data Element 55 ISO 8583, см.~\ref{sec:iso8583-data-elements}). EMV
использует подмножество ISO/IEC~8825-1: одно- и~двухбайтовые теги,
short и~long form длины, иногда вложенные \emph{constructed}-теги
(шаблоны/templates~--- конструктивные теги BER-TLV, содержащие
вложенные TLV; у~них бит~6 первого байта тега равен единице~--- например, FCI
Template \texttt{0x6F}, File Control Information). Минимальный парсер этого подмножества:

\codefile{Python}{samples/ch07-emv-tlv/tlv.py}

\section{Аутентификация данных}\label{sec:emv-data-auth}

Получив AIP и~AFL, терминал знает, какие файлы читать и~какой метод
аутентификации данных поддерживает карта. Аутентификация подтверждает,
что данные, записанные эмитентом при персонализации (этап выпуска карты,
на~котором эмитент через бюро персонализации загружает в~чип данные
приложения, ключи и~параметры), не были подменены после. EMV определяет
три метода аутентификации, расположенные по~возрастанию надёжности:
SDA, DDA и~CDA~\cite{emvco-book2}.

\subsection*{SDA: статическая аутентификация данных}

SDA (Static Data Authentication)~--- самый простой метод. При выпуске
карты эмитент формирует подпись RSA над определённым набором данных
карты, и~эта подпись записывается на~чип. При каждой транзакции
терминал считывает подписанные данные и~подпись, а~затем проверяет
её по~цепочке сертификатов: корневой сертификат карточной сети
заверяет сертификат эмитента, тот~--- подпись данных карты.

SDA подтверждает только, что данные не менялись после персонализации.
Подпись статична, и~если скопировать все данные с~чипа вместе с~подписью
на~другой чип, терминал подмену не заметит~--- подпись по-прежнему
валидна. SDA защищает от~модификации данных, но~не от~клонирования
чипа. В~современных профилях SDA вытеснен динамическими методами
и~остаётся в~обороте лишь как историческая ступень EMV.

\subsection*{DDA: динамическая аутентификация данных}

DDA (Dynamic Data Authentication) решает проблему клонирования: помимо
статической подписи эмитента, карта содержит собственную пару ключей
RSA. Приватный ключ хранится в~защищённой памяти чипа, а~открытый
ключ заверен сертификатом эмитента. Корень цепочки терминал выбирает
по~паре RID + \emph{CA Public Key Index} (тег \texttt{0x9F22}): для каждой
платёжной сети на~терминале хранится свой набор корневых ключей, и~индекс
указывает, какой именно ключ использовать при проверке сертификата
эмитента~\cite{emv-tag-9F22}. При каждой транзакции терминал отправляет
карте непредсказуемое число (unpredictable number, тег~\texttt{0x9F37})
в~составе DDOL (Dynamic Data Authentication Data Object List,
тег~\texttt{0x9F49})~--- списка тегов, значения которых карта ожидает
в~теле команды \texttt{INTERNAL AUTHENTICATE}. В~ответ карта формирует
\emph{Signed Dynamic Application Data} (тег~\texttt{0x9F4B})~--- подпись
над Dynamic Application Data, куда входят \emph{ICC Dynamic Number}
(тег~\texttt{0x9F4C}, вырабатываемый чипом для каждой транзакции)
и~значения тегов из~DDOL~\cite{emv-tag-9F4C}. Терминал проверяет эту
подпись, используя открытый ключ карты из~сертификата.

\begin{figure}[tbp]
  \centering
  \includefiguresvg[width=\linewidth,height=.7\textheight]{assets/figures/ch07-emv/emv-dda-sequence}
  \caption{Контактный сценарий DDA с~командой \texttt{INTERNAL AUTHENTICATE}.}\label{fig:emv-dda-sequence}
\end{figure}

\noindent Схема выше отражает контактный профиль с~командой
\texttt{INTERNAL AUTHENTICATE}. В~бесконтактных ядрах (kernel~---
реализация EMV-сценария на~стороне терминала под конкретную
платёжную сеть, подробно в~разделе~\ref{sec:contactless-emv})
SDAD формируется уже в~ответе на~\texttt{GET PROCESSING OPTIONS}
и~окончательно закрепляется командой \texttt{GENERATE AC}~---
APDU-командой EMV, по~которой карта вырабатывает финальную
application cryptogram (см.~раздел~\ref{sec:emv-cryptogram}). Отдельный
шаг \texttt{INTERNAL AUTHENTICATE} (APDU-команда EMV для запроса
динамической аутентификации карты) здесь не выполняется.
Бесконтактное ядро Visa~--- Kernel~3 в~EMVCo-нумерации; вариант
DDA в~нём называется fDDA (fast DDA).

Приватный ключ карты хранится в~области памяти чипа,
физически недоступной для чтения извне; он используется только
внутри криптосопроцессора и~никогда не покидает микросхему.
Скопируй злоумышленник все читаемые данные с~чипа~--- без
приватного ключа он всё равно не подпишет случайное число,
которое предъявит терминал в~следующей транзакции.
Криптографические примитивы за~подписями RSA разбираются
в~разделе~\ref{sec:crypto-asymmetric}~\cite{emvco-book2}.

\subsection*{CDA: комбинированная аутентификация}

CDA (Combined DDA / Application Cryptogram Generation) объединяет
динамическую аутентификацию данных и~генерацию криптограммы в~один
шаг, при котором карта подписывает своим приватным ключом не только
непредсказуемое число терминала, но и~саму криптограмму транзакции~---
ARQC (Authorization Request Cryptogram), TC (Transaction Certificate)
или AAC (Application Authentication Cryptogram); подробно эти три
типа криптограмм разобраны в~разделе~\ref{sec:emv-cryptogram}. Криптограмма тем самым привязана
к~конкретному чипу и~не может быть подменена в~процессе передачи.
Там, где карта и~терминал поддерживают CDA, аутентификация
сцеплена с~финальным решением о~транзакции заметно прочнее,
чем в~SDA или DDA~\cite{emvco-book2}.

\subsection*{Слабости DDA на~стороне терминала}

DDA устойчив ровно настолько, насколько случаен UN со~стороны
терминала. В~работе \enquote{Chip and Skim} (IEEE~S\&P, 2014) группа
из~Кембриджа показала, что эта посылка нарушалась
в~двух местах сразу~\cite{bond-chipandskim-2014}. Часть реальных
банкоматов и~POS формировала UN счётчиком, временной меткой или
самописным PRNG (pseudorandom number generator, программный
псевдослучайный генератор): предсказывая будущий UN, атакующий
заранее снимает с~карты криптограммы под нужные значения
и~проигрывает их позже~--- для эмитента это неотличимо
от~клонирования. Параллельно работает протокольная слабость:
UN рождается на~терминале, а~проверяется у~эмитента, поэтому
вредоносное ПО в~POS или MITM (man-in-the-middle, посредник
на~канале) на~участке \enquote{терминал~--- эквайер} может
подменить легитимный UN на~тот, под который криптограмма
у~атакующего уже есть.

Уязвимыми оказались терминалы крупных производителей. Аппаратный
TRNG чипа (см.~раздел~\ref{sec:emv-chip-internals}) защищает только
сторону карты, а~не сторону терминала~--- настоящая случайность UN должна
рождаться там, где он генерируется. Криптограмма поверх
скомпрометированного nonce уже не помогает.

\section{Верификация держателя}\label{sec:emv-cvm}

Аутентификация данных подтверждает, что карта подлинная. Но подлинная
карта может оказаться в~чужих руках~--- украденная, найденная, забытая
на~прилавке. Верификация держателя (Cardholder Verification) отвечает
на~другой вопрос: транзакцию инициирует законный владелец карты или нет.

EMV не навязывает единственного способа проверки. Карта хранит
упорядоченный список методов верификации держателя~--- CVM List
(Cardholder Verification Method List; CVM~--- Cardholder Verification
Method, метод верификации держателя)~\cite{emvco-book3}. Во~время транзакции
терминал идёт по~списку сверху вниз и~для каждого метода проверяет
применимость: есть ли клавиатура для PIN, превышена ли нужная сумма,
допускает ли конфигурация подпись или локальную проверку на~устройстве.
Выбирается первый применимый метод.

Основные семейства методов верификации.

\begin{itemize}
  \item \textbf{Offline PIN, открытый текст}~--- держатель вводит
    PIN на~клавиатуре терминала, терминал передаёт его чипу командой
    VERIFY, и~чип сравнивает введённое значение с~эталоном,
    хранящимся в~защищённой памяти. PIN никогда не покидает
    периметр \enquote{терминал + чип}. Open PIN~--- исторический вариант,
    который постепенно выводится из~сценариев без участия кассира~\cite{emvco-newsletter-q1-2025,visa-tadg}.
  \item \textbf{Offline PIN, зашифрованный}~--- аналогичен открытому,
    но PIN передаётся чипу в~зашифрованном виде с~использованием
    открытого ключа карты, что защищает от~перехвата на~интерфейсе
    терминал--чип.
  \item \textbf{Online PIN}~--- PIN вводится на~клавиатуре терминала,
    шифруется в~PIN-блок (о~форматах PIN-блока см.\
    в~разделе~\ref{sec:crypto-pin-block}) и~отправляется в~авторизационном
    сообщении эмитенту через DE~52. Чип в~проверке не участвует;
    верификацию выполняет хост эмитента.
  \item \textbf{Подпись}~--- кассир визуально сверяет подпись на~чеке
    с~образцом на~обратной стороне карты. Метод архаичный и~ненадёжный,
    но по-прежнему встречается в~некоторых конфигурациях.
  \item \textbf{No CVM}~--- верификация не требуется; обычно это мелкие
    суммы или определённые категории ТСП, например вендинговые
    автоматы.
  \item \textbf{Проверка на~устройстве}~--- биометрия или пароль
    на~мобильном устройстве вместо PIN на~терминале; этот метод
    применяется при оплате телефоном и~подробно разбирается
    в~разделе~\ref{sec:contactless-mobile}.
\end{itemize}

Точка доверия в~каждом случае разная~--- внутри чипа, у~эмитента,
у~кассира или на~устройстве держателя.

Если ни один метод из~CVM List не применим~--- скажем, у~терминала
нет клавиатуры, а~карта требует PIN,~--- результат фиксируется как
\enquote{CVM failed}. Соответствующий бит попадает в~TVR (Terminal
Verification Results, подробно~--- в~разделе~\ref{sec:emv-tvr})
и~влияет на~решение о~типе криптограммы на~последнем шаге сценария.

Исход верификации терминал записывает в~тег \texttt{0x9F34}
(CVM Results)~--- три байта: метод, условие применения, результат.
\texttt{0x9F34} уходит в~DE~55 вместе с~криптограммой. При отказе
эмитента \texttt{55} (Incorrect PIN) тег~--- первая точка расследования:
если он указывает на~подпись или No CVM, дело не в~PIN-блоке,
а~в~рассогласовании конфигурации терминала с~ожиданиями эмитента.

\section{TVR и~TSI: битовые поля результатов проверок}\label{sec:emv-tvr}

TVR (Terminal Verification Results)~--- пятибайтовое поле, в~которое
терминал складывает результаты проверок перед финальным решением.
Оно передаётся в~DE~55 под тегом \texttt{0x95}~\cite{emvco-book3}.

Каждый байт TVR отвечает за~свою группу проверок; нумерация идёт
от~старшего бита к~младшему внутри каждого байта (бит~8~--- MSB).

Установленная единица для большинства битов означает \enquote{шаг
не~прошёл} или \enquote{сработало нештатно}; в~логе отказанной
транзакции по~выставленным битам сразу видна причина отказа.
Расшифровка по~битам всех пяти байт~--- в~табл.~\ref{tab:tvr-bits};
жёлтая ячейка~--- бит~4 байта~3, установленный в~разобранном ниже
примере \texttt{TVR = 0x0000080000}.

Байты соответствуют пяти группам проверок.
\textbf{Байт~1}~--- offline data authentication (раздел~\ref{sec:emv-data-auth}): выполнялась~ли ODA и~какой из~методов SDA\slash DDA\slash CDA не~прошёл.
\textbf{Байт~2}~--- управление выбором приложения на~терминале: совпадение версии с~конфигурацией, срок действия, активация, разрешённые услуги, признак \enquote{новая карта}; даты~--- теги \texttt{5F24} (\textit{Application Expiration Date}) и~\texttt{5F25} (\textit{Application Effective Date})~--- читаются по~AFL уже после GPO, на~шаге PPSE FCI их ещё нет.
\textbf{Байт~3}~--- итог CVM (раздел~\ref{sec:emv-cvm}): выбранный из~CVM List метод, доступность клавиатуры, счётчик попыток PIN, факт ввода PIN.
\textbf{Байт~4}~--- риск-менеджмент терминала (раздел~\ref{sec:emv-risk-mgmt}): floor limit, лимит подряд офлайн-операций, случайный отбор, принудительный онлайн от~продавца.
\textbf{Байт~5}~--- состояние после ARQC и~онлайн-ответа эмитента (раздел~\ref{sec:emv-cryptogram}): сошёлся~ли ARPC, отработали~ли issuer scripts до и~после финальной криптограммы, использовался~ли TDOL по~умолчанию.

\begin{table}[H]
  \caption{TVR: расшифровка битов всех пяти байтов; жёлтым~--- бит~4 байта~3 примера \texttt{TVR = 0x0000080000}. Маска бита: \texttt{0x80}~--- бит~8, далее \texttt{0x40}, \texttt{0x20}, \texttt{0x10}, \texttt{0x08}, \texttt{0x04}, \texttt{0x02}, \texttt{0x01}. Имена флагов даны в~сокращённой форме EMVCo; полные формулировки~--- в~EMV Book~3, Annex~C5~\cite{emvco-book3}.}\label{tab:tvr-bits}
  \begingroup
  \scriptsize
  \setlength{\tabcolsep}{3pt}
  \renewcommand{\arraystretch}{1.2}
  \begin{tabularx}{\textwidth}{@{}P{1.1em} Y Y Y Y Y@{}}
    \toprule
    \headrow
    \thd{Бит} & \thd{Б1: ODA} & \thd{Б2: terminal} & \thd{Б3: CVM} & \thd{Б4: risk} & \thd{Б5: issuer auth} \\
    \midrule
    8 & ODA not performed & App version mismatch & CV not successful & Floor limit exceeded & Default TDOL used \\
    7 & SDA failed & Expired application & Unrecognised CVM & LCOL exceeded & Issuer auth failed \\
    6 & ICC data missing & Not yet effective & PIN Try Limit exceeded & UCOL exceeded & Script fail pre-GAC \\
    5 & Exception file & Service not allowed & No PIN pad & Random online & Script fail post-GAC \\
    4 & DDA failed & New card & \cellcolor{Warn!18}PIN not entered & Merchant forced online & --- \\
    3 & CDA failed & --- & Online PIN entered & --- & --- \\
    \bottomrule
  \end{tabularx}%
  \endgroup
\end{table}

Диагностический разбор примера \texttt{TVR = 0x0000080000}:
установленный бит~4 байта~3 означает, что карта потребовала PIN при
работающей клавиатуре, но держатель нажал \enquote{Отмена} или кассир
пропустил шаг. В~сочетании с~IAC Online, требующим онлайн-авторизацию
при неуспешной верификации, транзакция уходит онлайн и~получает отказ
эмитента \texttt{55} (Incorrect PIN)~--- хотя PIN физически не запрашивался.
Готовые онлайн-разборщики EMV-тегов (например, emvdecoder.com~\cite{emvdecoder})
декодируют TVR, AIP, TSI и~CVM List на~стороне клиента.

\subsection*{TSI: что было выполнено}\label{sec:emv-tsi}

Парный к~TVR тег~--- TSI (Transaction Status Information,
\texttt{0x9B}). Те же 2 байта битовых флагов, но смысл противоположный:
TVR фиксирует \emph{провалы и факты} проверок (\enquote{что не прошло}),
TSI~--- \emph{что вообще было выполнено} в~ходе
транзакции~\cite{emvco-book3}. Шесть определённых битов лежат в~байте~1
(табл.~\ref{tab:tsi-bits}); биты~2 и~1 байта~1 и~весь байт~2
зарезервированы (RFU). На~рис.~\ref{fig:tsi-anatomy} жёлтым отмечен
пример \texttt{TSI = 0xE800}~--- офлайн-транзакция, в~которой выполнены
ODA, cardholder verification, card risk management и~terminal risk
management; issuer authentication и~script processing не запрашивались,
потому что онлайн-обращения не было.

\begin{figure}[tbp]
  \centering
  \includefiguresvg[width=\linewidth]{assets/figures/ch07-emv/tsi-anatomy}
  \caption{Раскладка TSI: 2~байта, MSB слева; жёлтым~--- биты примера \texttt{TSI = 0xE800}.}\label{fig:tsi-anatomy}
\end{figure}

\begin{table}[H]
  \begingroup
  \small
  \setlength{\tabcolsep}{4pt}
  \renewcommand{\arraystretch}{1.2}
  \begin{tabularx}{\textwidth}{@{}P{1.4em} P{3em} Y@{}}
    \toprule
    \headrow
    \thd{Бит} & \thd{Маска} & \thd{Назначение} \\
    \midrule
    8 & \texttt{0x80} & Offline data authentication was performed \\
    7 & \texttt{0x40} & Cardholder verification was performed \\
    6 & \texttt{0x20} & Card risk management was performed \\
    5 & \texttt{0x10} & Issuer authentication was performed \\
    4 & \texttt{0x08} & Terminal risk management was performed \\
    3 & \texttt{0x04} & Script processing was performed \\
    \bottomrule
  \end{tabularx}%
  \endgroup
  \caption{Биты TSI (байт~1)~\cite{emvco-book3}.}\label{tab:tsi-bits}
\end{table}

\section{Риск-менеджмент терминала и~карты}\label{sec:emv-risk-mgmt}

После аутентификации данных и~верификации держателя остаётся
вопрос: одобрить транзакцию офлайн, без обращения
к~эмитенту, или отправить запрос онлайн. EMV отвечает на~него
двухуровневым риск-менеджментом~--- на~стороне терминала
и~на~стороне карты. Оба уровня записывают свои результаты
в~битовое поле TVR (Terminal Verification Results), из~которого
затем складывается итоговое решение~\cite{emvco-book3}.

\textbf{Лимиты терминала} смотрят на~текущую операцию~--- пороги
по~сумме и~правила случайного отбора; их выставляет эквайер
по~правилам схемы. \textbf{Прикладные лимиты карты} следят
за~историей самой карты~--- сколько подряд операций прошло
офлайн и~на~какую общую сумму; их записывает эмитент при
персонализации. Для бесконтактного канала действует ещё третий
набор~--- пороги CVM и~DRL (Dynamic Reader Limits, динамические
лимиты ридера), см.~\ref{sec:contactless-emv}.

\subsection*{Риск-менеджмент терминала}

Первая линия~--- \textbf{floor limit}, порог автономного одобрения
по~сумме. Эквайер выставляет его по~правилам схемы; у~схемы
и~эмитента могут быть собственные значения, и~терминал применяет
минимум. Сумма ниже порога допускает офлайн-одобрение; выше~--- терминал
обязан уйти онлайн и~поднимает в~TVR флаг \enquote{Transaction exceeds
floor limit} (байт~4, бит~8; см.~\ref{sec:emv-tvr}). Сегодня
у~большинства POS-терминалов floor limit равен нулю, и~любая
операция уходит на~эмитента (об~исходах
см.~\ref{sec:decline-offline-rc}).

Вторая линия~--- \textbf{случайный отбор}. Эмитент периодически
\enquote{видит} даже те операции, которые по~сумме могли бы пройти
офлайн. Ниже порога Threshold Value for Biased Random Selection
терминал отправляет онлайн фиксированную долю Target Percentage
for Random Selection; между этим порогом и~floor limit вероятность
линейно растёт до~Maximum Target Percentage for Biased Random
Selection по~мере приближения суммы к~floor limit
(EMV Book~3, \S10.6.2; рис.~\ref{fig:emv-random-selection})~\cite{emvco-book3}.

\begin{figure}[tbp]
  \centering
  \includefiguresvg[width=.7\textwidth]{assets/figures/ch07-emv/emv-random-selection}
  \caption{Вероятность ухода операции онлайн при случайном отборе.}\label{fig:emv-random-selection}
\end{figure}

Третья~--- контроль последовательных офлайн-операций по~лимитам,
записанным в~данных карты: \textbf{Lower\slash Upper Consecutive
Offline Limit}. Если карта слишком давно не выходила онлайн, терминал
больше не вправе одобрять операцию офлайн.

\subsection*{Риск-менеджмент карты}

Карта ведёт собственный учёт~--- счётчик транзакций ATC (Application
Transaction Counter, тег \texttt{0x9F36}), кумулятивные суммы
офлайн-транзакций, количество последовательных офлайн-операций. Текущие
значения карта сравнивает с~лимитами, установленными эмитентом
при персонализации, и~при исчерпании лимитов сама уводит транзакцию
онлайн или отклоняет её~--- независимо от~того, что решил терминал.

\subsection*{TAC и~IAC: матрица решений}

Результаты всех проверок собираются в~TVR~--- пятибайтовое поле, где
каждый бит отражает итог конкретной проверки; подробный разбор TVR
см.\ в~разделе~\ref{sec:emv-tvr}. Чтобы превратить набор признаков
в~исход транзакции, терминал сопоставляет TVR с~двумя наборами масок.
TAC (Terminal Action Codes) настраивает эквайер~--- какие сочетания
признаков ведут к~онлайн-авторизации, отклонению или офлайн-допуску.
IAC (Issuer Action Codes) записывает эмитент при персонализации; это
его собственная политика на~тот же набор ситуаций.

Терминал делает побитовое AND между TVR и~соответствующими масками
TAC/IAC. Ненулевой результат для маски \textit{Denial} ведёт к~отклонению
и~AAC. Ненулевой результат для маски \textit{Online} уводит операцию в~онлайн
и~приводит к~ARQC. Если ни одно условие не срабатывает, карта может
одобрить транзакцию офлайн и~выдать TC. Современные конфигурации
из-за низких офлайн-порогов почти всегда идут онлайн, но офлайн-ветвь
сохранена в~стандарте, и~флаги TVR её по-прежнему обслуживают.

\section{Криптограмма: ARQC, TC, AAC}\label{sec:emv-cryptogram}

Криптограмма завершает транзакцию EMV. Восьмибайтовое значение,
вычисленное чипом по~секретному ключу, одновременно
фиксирует решение карты и~доказывает подлинность чипа. По~итогам
предыдущих фаз (аутентификации данных, верификации держателя
и~риск-менеджмента) карта генерирует одну из~трёх
криптограмм~\cite{emvco-book3}.

\textbf{ARQC (Authorization Request Cryptogram)}~--- криптограмма,
с~которой карта просит онлайн-авторизацию у~эмитента; самый
распространённый исход контактной транзакции EMV. По~настройкам
IAC, превышению лимитов или требованию терминала карта определяет,
что одобрение офлайн невозможно, и~формирует криптограмму для эмитента.

\textbf{TC (Transaction Certificate)}~--- \enquote{карта одобряет
транзакцию офлайн}. Эмитент в~решении не участвует; TC уходит
в~клиринг как доказательство, что транзакция была одобрена чипом.

\textbf{AAC (Application Authentication Cryptogram)}~--- \enquote{карта
отклоняет транзакцию}. Терминал должен остановиться.

Все три исхода представлены как разные типы application
cryptogram. Тип дополнительно кодируется в~теге \texttt{0x9F27}
(Cryptogram Information Data)~\cite{emvco-book3}. В~публичном
глоссарии Visa TADG ARQC, TC и~AAC описаны как результат
обработки данных карты, терминала и~транзакции ключом TDEA или
AES~\cite{visa-tadg}; детали иерархии ключей и~деривации~---
в~разделе~\ref{sec:crypto-key-hierarchy}. В~авторизационном
сообщении криптограмма передаётся внутри DE~55 под тегом
\texttt{0x9F26}; см.~раздел~\ref{sec:iso8583-data-elements}.

\begin{figure}[tbp]
  \centering
  \includefiguresvg[width=\linewidth]{assets/figures/ch07-emv/emv-cryptogram-decision}
  \caption{Принятие решения о~типе криптограммы в~EMV.}\label{fig:emv-cryptogram-decision}
\end{figure}

\subsection*{CVN: какой именно вариант криптограммы}

Чтобы хост эмитента воспроизвёл ARQC, ему нужен один параметр сверх
данных транзакции~--- \textbf{CVN (Cryptogram Version Number)}.
Он передаётся внутри IAD (Issuer Application Data, тег \texttt{0x9F10})
и~говорит хосту: каким способом карта прошла иерархию ключей
от~мастер-ключа к~ключу карты и~далее к~сессионному ключу
(см.~раздел~\ref{sec:crypto-key-hierarchy}) и~какой формат MAC она применила.
CVN~10 и~18 на~одних и~тех же данных карты, ATC, UN и~IMK дают разные
восьмибайтовые значения, и~типичная причина расследования несошедшейся
криптограммы~--- профиль персонализации карты под одну версию,
а~настройки хоста~--- под другую.

В~семействе Visa~VIS три CVN встречаются чаще остальных~\cite{visa-tadg}:

\begin{itemize}
  \item \textbf{CVN~10}~--- ключ карты (UDK, Unique Derivation Key)
    выводится из~IMK по~PAN и~PSN (PAN~Sequence~Number), и~UDK
    используется как сессионный ключ напрямую, без диверсификации
    по~ATC. MAC считается по~CDOL1 (Card Risk Management Data Object
    List~1; список тегов, которые карта запрашивает у~терминала
    при первом \texttt{GENERATE~AC}), ARPC формируется методом~1
    (8~байт, XOR криптограммы с~ARC~--- Authorization Response
    Code, кодом ответа эмитента~--- и~шифрование сессионным ключом).
  \item \textbf{CVN~18}~--- поверх UDK выводится сессионный ключ
    по~ATC, как в~разделе~\ref{sec:crypto-key-hierarchy}. ARPC формируется
    методом~2 (4~байта, MAC поверх ARQC и~CSU~--- Card~Status~Update,
    инструкций эмитента к~карте).
  \item \textbf{CVN~22}~--- то же дерево, что у~CVN~18, но AES
    вместо 3DES; вводится в~ходе общей миграции платёжной
    криптографии с~TDEA на~AES (см.~раздел~\ref{sec:crypto-symmetric}).
\end{itemize}

Mastercard в~профиле M/Chip использует собственный CVN с~ARPC
по~методу~2; внутренние детали закрыты~\cite{mc-mchip}.

\subsection*{ARPC: ответ эмитента}

Когда эмитент получает авторизационный запрос с~ARQC, он проверяет
криптограмму своим ключевым материалом и~данными операции.
Совпадение подтверждает, что карта аутентична, а~защищённые
криптограммой данные карты и~терминала не были подменены
в~пути~\cite{visa-tadg}. Затем эмитент одобряет или отклоняет операцию
и~формирует ARPC (Authorization Response Cryptogram)~--- ответную
криптограмму issuer authentication. Карта проверяет ARPC, чтобы
убедиться, что ответ действительно пришёл от~легитимного эмитента
и~не был изменён в~пути~\cite{visa-tadg}.

Терминал передаёт ARPC чипу. Если ARPC валиден и~эмитент одобрил
транзакцию, карта перегенерирует криптограмму, теперь TC,
подтверждая итог. Если ARPC невалиден или эмитент отклонил, карта
генерирует AAC.

\subsection*{Сценарии эмитента (issuer scripting)}

Вместе с~ARPC эмитент может отправить чипу сценарии управления
(issuer scripts). Они передаются в~тегах EMV \texttt{0x71}
(Script Template 1, выполняется до~финальной криптограммы)
или \texttt{0x72} (Script Template 2, выполняется после). Через них
эмитент обновляет параметры карты~--- сбрасывает офлайн-счётчик
попыток PIN, меняет офлайн-лимиты, блокирует приложение или всю карту.
Прямой канал к~карте не нужен. Достаточно, чтобы держатель совершил
любую транзакцию через любой терминал~\cite{emvco-book3}.

\section{Возврат к~магнитной полосе (fallback)}\label{sec:emv-fallback}

Не каждая чиповая транзакция доходит до~конца на~чипе. Контакт
между терминалом и~площадкой бывает нестабилен (карта изношена,
загрязнена, терминал повреждён), приложение на~чипе не отвечает,
обмен APDU ломается на~полпути. В~таких случаях терминал
выполняет \textit{fallback}~--- переключается на~считывание магнитной полосы
и~проводит транзакцию как обычную операцию по~полосе~\cite{emvco-book4}.

Fallback бывает двух типов. Технический: терминал не смог прочитать
чип и~после нескольких неудачных попыток сам перешёл на~магнитную
полосу. Авторизационное сообщение несёт специальное значение
POS Entry Mode в~DE~22, и~эмитент понимает: чип не прочитался,
операция идёт по~полосе. Обходной: обход инициирует сама торговая
точка. Кассир проводит карту через магнитный считыватель, минуя чип.
Причиной бывает спешка, неисправный терминал или злонамеренное
поведение; именно этот сценарий карточные сети ограничивают жёстче
всего.

\textit{Fallback} возвращает транзакцию к~статическим данным магнитной полосы
вместо динамической криптограммы чипа. Карточные сети рассматривают
его как исключение~--- он ухудшает риск-профиль и~может перераспределять
ответственность. В~публичном Visa TADG жёстче: транзакции с~\textit{fallback}
должны авторизовываться онлайн, и~порог офлайн-авторизации к~ним
не применяется~\cite{visa-tadg}. Пороги мониторинга,
штрафы и~нюансы правил привязаны к~каждой схеме отдельно
и~разбираются в~главах~\ref{ch:visa} и~\ref{ch:mastercard}.

Дополнительной линией защиты в~сценарии \textit{fallback} остаётся
iCVV (см.~\ref{sec:card-data-cvv}). Он отсекает попытку выдать
скопированные с~чипа данные за~магнитную полосу, но не помогает
против обычного скимминга самой полосы. Конечная цель миграции
на~EMV~--- полный отказ от~магнитной полосы~\cite{visa-core-rules}.

\section{Скорость обмена: EFA и~Fast/Quick Chip}\label{sec:emv-fast-flows}

Заметная пауза на~кассе при оплате контактным чипом складывается
из~длины APDU-цепочки к~карте и~накладных расходов протоколов
T=0/T=1.

Сверх этого ряд протоколов между POS и~ПЦ эквайера
требует второй сессии связи после онлайн-ответа эмитента. EMVCo описывает её как online
completion (завершение онлайн-транзакции на~чипе после ARPC) /
second GENERATE~AC (повторный вызов команды \texttt{GENERATE AC}
после онлайн-ответа эмитента, на~котором карта генерирует финальную
криптограмму TC или AAC); в~процессорной практике за~этим
шагом закрепилось индустриальное название EMV Final Advice (EFA),
и~обычно он несётся в~ISO~8583 как пара MTI (Message Type Indicator,
тип сообщения ISO~8583, см.~\ref{sec:iso8583-data-elements})~0220
(advice от~терминала)
и~0230 (ответ хоста). В~ней терминал передаёт на~хост результат
обработки картой ARPC, финальную криптограмму TC или AAC и~данные,
нужные для формирования клиринговой записи~\cite{emvco-book4}.
POS дважды открывает соединение: сначала чтобы получить ответ
эмитента, затем чтобы сообщить итог транзакции.

Чтобы убрать издержки, схемы выпустили облегчённые сценарии
под общим именем \enquote{Fast} или \enquote{Quick}: Mastercard
M/Chip Fast, Visa Quick Chip, AmEx Quick Chip. Идея общая~---
завершить взаимодействие с~картой сразу после первого
GENERATE~AC, не дожидаясь ответа эмитента и~не запрашивая вторую
криптограмму. Карта возвращает либо TC, либо ARQC; авторизационный
обмен с~ПЦ становится опциональным и~асинхронным относительно
держателя, который уже забрал чек и~ушёл~\cite{visa-tadg,mc-tpr}.

В~российской региональной инфраструктуре Fast/Quick Chip широкого
распространения не получил. Задачу скорости решили иначе~--- быстрым
переходом на~бесконтакт, у~которого радиоинтерфейс и~упрощённое
ядро дают сравнимое время обслуживания без переделки логики
на~стороне эмитента; см.~главу~\ref{ch:contactless}.
